\chapter{Barnløse Joakim og Anna}

I dagene da Israels folk fremdeles levde etter fedrenes skikker,
fantes det en mann ved navn Joakim. Han var rik på jord,
buskap og eiendommer, men det som gjorde størst inntrykk
på de som kjente ham, var ikke rikdommen hans. Det var hans fromhet.

Joakim trodde at alt han eide, kom fra Gud. Derfor holdt han ikke
hardt fast på sine rikdommer. Når høytidene kom og folket samlet
seg for å bære fram offergaver, gav han mer enn loven krevde.
Han delte gavene sine i to deler.
Den ene delen gav han til de fattige, de foreldreløse og de nødlidende.
Den andre delen viet han som forsoningsoffer til Herren.
Ofte sa han: «Hvis Gud har velsignet meg mer enn andre,
skal også andre få del i velsignelsen. Og det jeg gir til Herren,
gir jeg med takknemlighet og ydmykhet.»  

Likevel var det én sorg som aldri forlot ham.
Han og hans hustru Anna hadde levd sammen i mange år.
De elsket hverandre og hadde aldri mistet håpet om å få barn.
Men årene gikk. Hver vår kom med nytt liv i markene.
Hver høst så de barn løpe mellom oliventrærne og vinrankene.
Hver gang et nytt barn ble født i landsbyen,
gledet de seg på foreldrenes vegne.
Men deres eget hus forble stille.
I den tiden var barnløshet mer enn en personlig sorg.
Mange mente at en stor familie var et tegn på Guds gunst.
Derfor kunne mennesker hviske bak ryggen på
dem som ikke hadde barn. Noen syntes synd på dem. Andre dømte dem.

Da den store løvhyttefesten  nærmet seg, dro Joakim til tempelet sammen
med de andre pilegrimene. Veiene var fulle av folk som sang
salmer og bar med seg offerdyr og gaver. Luften var fylt av forventning.
Joakim kom fram med sine offergaver, slik han alltid hadde gjort.
Men før han rakk å legge dem fram, trådte en mann ved navn
Rubel fram rett foran ham. Mannen løftet stemmen slik at alle andre skulle
kunne høre.
 «Hvorfor kommer du først? Du har ikke etterkommere i Israel.
Hvordan kan du regne deg blant de velsignede?»  
Ordene falt tungt.
Samtalene omkring dem stilnet. Noen vendte blikket bort. Andre stirret.
Joakim svarte ikke.
Han kjente varmen stige i ansiktet. Hjertet sank i brystet på ham.
Han hadde tålt mange år med sorg, men aldri tidligere hadde
den blitt stilt ut offentlig på denne måten.
Han snudde seg og gikk bort fra folkemengden.

Senere samme dag begynte han å lete i slektsbøkene og
fortellingene om Israels folk. Han gransket navnene på
patriarkene og de rettferdige som hadde levd før ham.
Kveld etter kveld satt han bøyd over rullene.
Han fant menn som hadde vært fattige og menn som hadde vært rike.
Han fant krigere, hyrder, dommere og profeter.
Men alle hadde de fått barn.
Til slutt kom han til fortellingen om Abraham.
braham hadde også ventet lenge. Han hadde kjent den samme sorgen.
Han hadde blitt gammel uten å se løftet oppfylt.
Likevel hadde Gud til slutt gitt ham Isak.

Joakim lukket øynene.
Hvis Gud hadde hørt Abraham, hvorfor var himmelen da så stille nå?
Sorgen vokste seg tung.

Da han vendte hjem, gikk han ikke inn til Anna.
Han kunne ikke få seg til å møte hennes blikk.
Ikke fordi han elsket henne mindre, men fordi sorgen
mellom dem var blitt så stor.
I stedet forlot han huset og drog ut mot ødemarken.
På  en steinete skråning ved en liten høyde vokste et gammelt oliventre. Det stod som om
det hadde voktet landskapet siden tidenes morgen. Den grove stammen var vridd og sprukket,
formet av utallige vintre med regn og somre under brennende sol.
Dype furer løp gjennom barken som rynkene i ansiktet til et menneske
som hadde levd et langt og begivenhetsrikt liv. 
Grenene strakte seg utover i alle retninger, krokete og seige,
som værbitte armer som nektet å gi etter for tidens tyngde.
Noen bøyde seg nesten ned mot bakken før de løftet seg igjen,
mens andre bar kroner av sølvgrønne blader som hvisket lavmælt når
den milde vinden gled gjennom dem. Mellom bladene hang små,
mørke oliven som stille vitner om at treet, til tross for sin høye alder,
fremdeles bar frukt. 
Røttene brøt gjennom den tørre jorden og klemte seg fast rundt steinene med en urokkelig styrke,
som om de omfavnet selve fjellet. Hver sprekk i stammen, hver avknekt gren og hver
knudrete utvekst fortalte om stormer som var overlevd, tørke som var utholdt,
og århundrer som var kommet og gått. I skyggen under kronen var luften svalere.
Der hvilte en stillhet som føltes gammel, nesten hellig, som om treet bar på minner
ingen mennesker lenger kunne huske. Det sto der uten hastverk, urørlig og tålmodig,
som en taus vokter over landet og over alle som gjennom generasjonene
hadde søkt ly under dets eldgamle grener.

Der, langt fra landsbyene og de travle veiene, like ved det eldgamle oliventreet
slo han opp et enkelt telt.


Om dagen brente solen over landet.
Om natten lå stjernene som tusen lys over himmelhvelvingen.
Joakim ble alene med sine tanker.

Han bestemte seg for å faste.
Dagene gikk. Så ble de til uker.
Han spiste ikke. Han drakk ikke annet
enn det aller nødvendigste for å overleve. Han viet tiden til bønn.

Gang på gang løftet han hendene mot himmelen og sa:
«Herre, mine fedres Gud, se til meg. Hør min bønn.
Jeg vil ikke vende tilbake før du viser meg barmhjertighet.»
Slik ble bønnene hans hans mat. Lengselen etter Gud ble hans drikke.
Og mens dagene gikk gjennom den stille ødemarken,
visste ikke Joakim at hans liv snart skulle forandres for alltid





Anna satt alene i huset sitt og stirret ut gjennom vinduet.

«Jeg sørger fordi jeg er blitt enke, og fordi jeg ikke har fått et barn.
Min ensomhet er dobbel; jeg har mistet min ektemann,
og jeg har ingen sønn eller datter som kan bære hans minne
videre eller gi trøst i mine eldre dager.», tenkte hun.

    Da dagen  for Herrens store høytid hadde kommet,
    kom tjenestekvinnen hennes, Juthine, inn til henne.
    «Hvor lenge vil du fortsette å å sørge?» spurte hun forsiktig.
    «Det er tid for fest og glede. Du kan ikke gjemme deg bort for alltid.»
   Men Anna ristet på hodet.

   «Hvordan kan jeg glede meg?» svarte hun. «Jeg er alene,
   og jeg har ingen barn. Hver dag minner meg om det jeg har mistet.»

   Juthine prøvde å oppmuntre henne. Hun ga Anna et vakkert
   hodebånd og sa   «Ta heller dette hodebåndet. Jeg kan ikke bære det selv.
   Det er prydet med et kongelig merke som passer
   bedre for en kvinne av din rang enn for n tjenestekvinne som meg.»

     Men Anna svarte henne:
    «Gå bort fra meg! Jeg vil ikke gjøre dette. Herren Gud
    har allerede latt meg erfare dyp ydmykelse og sorg.
    Kanskje har en bedrager gitt deg dette smykket,
    og kanskje er du kommet for å få meg til å ta del i en
    synd eller en ugjerning som ikke er min.»

    Da svarte Juthine, såret over å bli avvist:
    «Hvorfor skulle jeg ønske deg ulykke eller
    tale ondt om deg når du ikke engang vil høre på mine ord?
    Herren Gud har lukket ditt morsliv og nektet deg fruktbarhet,
    slik at du ikke har fått noen etterkommere i Israel.»

    Disse ordene traff Anna som et sverd i hjertet.
    Hun ble fylt av enda større sorg og fortvilelse.
    Likevel reiste hun seg. Hun tok av seg sørgeklærne
    som hadde vært hennes daglige drakt gjennom lang tid,
    vasket håret og renset seg, og deretter kledde hun seg i
    de vakre klærne hun hadde båret på sin bryllupsdag.

    Omkring den niende time på dagen gikk hun ned
    i hagen sin for å vandre alene. Hagen lå stille
    under et mildt lys, og der fikk hun øye på et staselig laurbærtre.
    Hun gikk bort til det og satte seg i skyggen under grenene.
    Luften under laurbærtreet bar en egen duft. De blanke bladene
    hadde samlet solens varme, og når vinden strøk gjennom grenene,
    spredte den en krydret aroma som blandet seg med lukten av
    jord og stein. Det luktet friskt og gammelt på samme tid ---
    som en plante som hadde stått der i mange år
    og bevart minnene fra alle årstidene den hadde sett.
    En svak bitterhet lå under den søte urteduften,
    som om treet bar på både styrke og sorg.
    Etter en stunds hvile løftet Anna sitt hjerte til Herren i bønn:    
    «Du mine fedres Gud, velsign meg og hør min bønn,
    slik du velsignet Sara vår mor og gav henne sønnen Isak.»


    
    Anna satt fortsatt i hagen under laurbærtreet.
    Hjertet hennes var tungt etter mange år med sorg,
    men hun hadde likevel løftet sin bønn mot Gud.
    Hun hadde ikke lenger flere ord igjen,
    bare håpet om at Han som hadde hørt de gamle fedres bønner,
    også kunne høre hennes.


    

Da skjedde noe uventet.
Luften rundt henne ble stille.
En mildt skinnende skikkelse viste seg nær henne.
Anna ble redd og reiste seg raskt.

«Anna, Anna.», sa engelen Gabriel med rolig stemme full av fred.
Hun så opp. «Ikke vær redd. Herren har hørt din bønn.»
Anna stod helt stille.

«Min bønn?» hvisket hun.

Gabriel svarte: «Ja. De tårene du har grått i det skjulte,
har ikke vært usett. Din sorg har nådd fram til Gud.
Du skal bli med barn, og du skal føde en sønn eller en datter.
Barnet som kommer fra deg, skal bli kjent langt utover dette landet.
Mange mennesker skal høre om det.»

Anna klarte nesten ikke å forstå ordene.
I så mange år hadde hun forberedt seg på å leve uten barn.
Hun hadde lært seg å bære skammen, stillheten og ensomheten.
Nå fikk hun høre at alt skulle forandre seg.

Hun falt ned på kne.
Tårene rant, men denne gangen var de ikke bare tårer av smerte.
De var tårer av glede.

«Så sant Herren min Gud lever», sa hun, «hvis du har gitt meg dette barnet,
skal jeg ikke holde det bare for meg selv.»
Hun løftet hendene.
«Enten barnet blir en gutt eller en jente, skal jeg gi det tilbake til Gud.
Hele livet skal barnet tjene Ham og tilhøre Ham.»

Engelen så på henne med et mildt blikk.
For Anna var dette ikke bare et ønske om å få et barn.
Det var et løfte.

Langt borte, på markene utenfor byen, var Annas ektemann,
Joakim, med flokkene sine.
Han var en rik mann, kjent for å være gavmild og rettferdig.
Men også han bar en sorg.
Han hadde i mange år bedt Gud om et barn.
En gang hadde han reist til tempelet med offergaver,
men blitt avvist fordi han og Anna var barnløse.
Dette hadde såret ham dypt.

Siden den dagen hadde han søkt ensomheten blant flokkene sine.
Der kunne han be uten at noen hørte.
Der kunne han gråte uten at noen så.
Men også til ham kom et budskap.

En engel viste seg for ham og sa: «Joakim.»
Han falt tilbake av frykt.

«Ikke vær redd», sa engelen.
«Herren Gud har hørt din bønn. Du har trofast båret sorgen din.
Nå skal du vende tilbake hjem.»

Joakim så opp.
«Hjem?» spurte han.

«Ja, din kone Anna har unnfanget et barn», sa engelen.
Ordene slo mot ham som en storm. Han klarte ikke å svare.
«Barnet du har ventet på, er gitt dere av Gud.»
Joakim reiste seg straks.

Han kalte på hyrdene som passet flokkene.
«Gjør klar de beste dyrene vi har», sa han.
Hyrdene samlet seg rundt ham.

«Hva skal vi bringe?»

Joakim svarte:
«Finn ti lam uten feil eller skade. De skal gis som gave til Herren.»
Han fortsatte:
«Finn også tolv unge okser til prestene og folkets ledere.
Og hundre geitebukker til folket, som et tegn på takknemlighet.»
Hyrdene så på hverandre.
De forstod at noe stort hadde skjedd.
Joakim, som i så mange år hadde båret sorg, hadde fått håpet tilbake.

Da Joakim nærmet seg huset deres, sto Anna ved porten.
Hun visste ennå ikke nøyaktig når han kom.
Men hjertet hennes kjente det.
Hun løftet blikket og så flokkene komme nedover veien.
Midt blant dyrene gikk Joakim, mannen hennes.
Den hun hadde delt livet med.
Den hun hadde sørget sammen med.
Et øyeblikk stod hun helt stille.
Så løp hun mot ham.

Joakim rakk knapt å åpne armene før hun kastet seg rundt halsen hans.
Begge gråt.
Ikke av sorg.
Men av lettelse.

«Nå vet jeg det», sa Anna med skjelvende stemme.
«Nå vet jeg at Herren virkelig har velsignet meg.»
Hun holdt rundt ham.
«Jeg var som en enke uten håp.»
Hun smilte gjennom tårene.
«Men nå er jeg ikke lenger forlatt.»

Hun la hånden på magen.
«Den som var uten barn, bærer nå et nytt liv.»
Joakim lukket øynene og takket Gud.
I mange år hadde huset deres vært fylt av stillhet.
Nå var stillheten brutt.
Et nytt kapittel hadde begynt.

Den første dagen etter at Joakim kom hjem, hvilte han i huset sitt,
men alt var annerledes.
Veggene var de samme.
Hagen var den samme.
Laurbærtreet stod fortsatt der.
Men sorgen som hadde bodd i huset så lenge, hadde fått vike.
For et løfte var blitt gitt, og et barn var på vei.


Dagen etter at Joakim hadde kommet tilbake til hjemmet sitt, våknet han tidlig.
Solen hadde ennå ikke rukket å stige høyt over åsene
da han gjorde seg klar til å gå til tempelet i Jerusalem.
Han hadde med seg gavene han hadde lovet Herren --- de beste
dyrene fra flokkene sine, valgt ut med omsorg og takknemlighet.
Men i hjertet bar han fortsatt et spørsmål.

I mange år hadde han båret skammen over å være barnløs.
Han hadde bedt. Han hadde ofret. Han hadde ventet.
Nå, etter engelens ord og Annas glede, våget han å håpe.
Men han ønsket en bekreftelse.
På vei til tempelet sa han stille til seg selv:
«Hvis Herren virkelig har tatt imot meg igjen,
hvis Han har tilgitt meg og vendt sitt ansikt mot meg,
da vil Han vise meg et tegn.»

Han visste at prestene bar en hellig plate på brystet
når de tjente foran alteret. For folket var den et tegn på Guds nærvær og dom.
«Når presten står framfor Herren», tenkte Joakim,
«skal jeg vite om min gave er mottatt.»

Tempelet strålte i morgensolen.
Mennesker kom fra alle kanter. Noen bar offergaver.
Andre kom for å be. Luften var fylt av lyden av bønner,
sang og skritt mot de store steinhellene.

Joakim gikk fram med gavene sine.
Han bøyde hodet og overlot dem til Herren.
Presten tok imot dem og gikk mot alteret.
Joakim fulgte nøye med. Han lette etter et tegn.
Et øyeblikk var hjertet hans fullt av frykt.
Men så skjedde det.
Han så ingen mørkhet.
Ingen tegn på skyld.
Ingen avvisning.
Bare fred.
En dyp ro fylte ham.

Joakim trakk pusten.
«Nå vet jeg det», hvisket han.
«Herren Gud har hørt meg. Han har tatt imot meg.»
Tårene kom til øynene hans.
Ikke tårer av sorg, men av lettelse.
«Mine feil er ikke lenger foran meg.
Min skam er ikke det siste ordet i livet mitt.»
Han bøyde hodet i takknemlighet.
Så forlot han tempelet.
Ikke lenger som en mann som ventet på dom.
Men som en mann som hadde fått håpet tilbake.
